% ============================================================
% Maxwell--SST main-canon patch block
% Target: SST_CANON-v0.8.35.tex
%
% Placement A:
% Replace the existing block beginning
%   \subsection{Ungapped Spectrum: Instability}
% through the end of the existing
%   \subsection{Consistency Condition}
% with the block below, then retain the subsequent
%   \subsection{Observable Consequences}.
% ============================================================
    \subsection{Ungapped Spectrum: Instability}

        If the internal spectrum is ungapped, arbitrarily low-energy
        excitations are kinematically available and no universal Boltzmann
        suppression follows from a gap alone.

        \textbf{[CRITICAL / CONDITIONAL].}
        An ungapped internal sector is generically incompatible with precision
        spectroscopy if its modes are appreciably coupled and equilibrate on the
        observation timescale.  Compatibility requires at least one independently
        demonstrated suppression mechanism, for example:
        \begin{itemize}
            \item $|\lambda_n|\ll1$ (extremely weak observable coupling);
            \item a sufficiently sparse or symmetry-forbidden active density of states;
            \item $\tau_{K,a}\gg t_{\rm obs}$ for the relevant interaction
            (dynamical freeze-out).
        \end{itemize}
        Thus ``ungapped'' does not by itself mean ``observably excited'', but an
        ungapped mode may not be deleted from the ledger without a coupling and
        timescale calculation.

    \subsection{Gapped Spectrum: Conditional Decoupling Mechanism}
        \label{subsec:canon_true_internal_gap}

        \textbf{[BRIDGE / DEFINITION].}
        A true excitation gap for knot class $K$ is an energy difference,
        \begin{equation}
            E_{m,n}-E_{0,n}
            \geq
            \Delta_K
            >0 .
            \label{eq:canon_true_excitation_gap}
        \end{equation}

        \textbf{[ORTHODOX INPUT / CONDITIONAL BRIDGE].}
        If the internal states are discrete and are in thermal equilibrium at a
        physically defined temperature $T_{\rm eff}$, their excited-state
        occupations are Boltzmann suppressed,
        \begin{equation}
            \nu_{m,n}
            \propto
            \exp\!\left(-\frac{E_{m,n}-E_{0,n}}{k_B T_{\rm eff}}\right),
            \label{eq:boltzmann_suppression}
        \end{equation}
        and therefore the leading contribution of a state separated by
        $\Delta_K$ is exponentially small when
        $\Delta_K/(k_BT_{\rm eff})\gg1$ \cite{PathriaBeale2011}.

        \textbf{[CRITICAL NOTE].}
        This statement is conditional on a genuine discrete energy gap and on
        equilibration.  It is not a theorem that every nonzero SST normal-mode
        frequency produces such a gap.

    \subsection{Physical Origin of a True Gap}

        \textbf{[SPECULATIVE / RESEARCH TARGET].}
        A positive $\Delta_K$ may arise only from a mechanism that excludes
        arbitrarily small excitations from the admissible state space, for
        example:
        \begin{itemize}
            \item a discrete occupancy or mode-selection rule;
            \item a topological barrier between admissible states;
            \item a nonlinear activation threshold;
            \item a finite-core transition with a derived minimum energy cost.
        \end{itemize}

        \textbf{[ORTHODOX INPUT / CANON GUARD].}
        Curvature stiffness, torsional rigidity, and a nonzero Kelvin frequency
        can generate a restoring force, but in a classical quadratic continuum
        they do not by themselves imply a discrete energy gap:
        \begin{equation}
            \omega_{K,a}>0
            \not\Rightarrow
            \Delta_{K,a}>0 .
            \label{eq:canon_gap_frequency_guard_local}
        \end{equation}

    \subsection{Dimensional Guard on Gap Estimates}

        \textbf{[DERIVED / CRITICAL CORRECTION].}
        No numerical internal-mode gap is canonized without a resolved energy
        functional.  In particular, an expression of the form
        \begin{equation}
            \frac{\Gamma^2}{\xi L}
        \end{equation}
        is not an energy by itself, because
        \begin{equation}
            \left[\frac{\Gamma^2}{\xi L}\right]
            =
            \frac{{\rm m^4\,s^{-2}}}{{\rm m^2}}
            =
            {\rm m^2\,s^{-2}} .
            \label{eq:canon_gap_dimensional_guard}
        \end{equation}
        A dimensionally valid gap must be computed as an actual energy
        difference of the same declared finite-core model,
        \begin{equation}
            \Delta_{K,a}
            =
            E_K[\mathcal Q_{K,a}^{\rm exc}]
            -
            E_K[\mathcal Q_K^{(0)}],
            \label{eq:canon_gap_energy_difference}
        \end{equation}
        or through another independently derived map with units of joules.
        The earlier order-of-magnitude claim
        $\Delta_K=\mathcal O(10^2\!-
        10^3\,\mathrm{eV})$ is therefore not retained as a derived Canon result.

    \subsection{Consistency Condition}

        \textbf{[DERIVED / CONSISTENCY GATE].}
        Compatibility with low-energy spectroscopy can be obtained by any
        independently established combination of:
        \begin{align}
            \frac{\Delta_{K,a}}{k_BT_{\rm eff}} &\gg 1
            &&\text{(gapped and equilibrated)},
            \label{eq:canon_gap_freeze_condition}\\
            |\lambda_{n;K,a}| &\ll 1
            &&\text{(weakly coupled)},
            \label{eq:canon_weak_mode_coupling}\\
            \tau_{K,a} &\gg t_{\rm obs}
            &&\text{(dynamically frozen)}.
            \label{eq:canon_mode_frozen_condition}
        \end{align}
        A numerical gap is not required if the mode is demonstrably decoupled or
        cannot equilibrate on the relevant observation timescale; conversely, a
        nominal gap does not protect the theory if strong driven transitions
        remain kinematically accessible.


    \subsection{Maxwell--SST Internal-Mode Thermodynamic Guard}
        \label{subsec:maxwell_sst_internal_mode_guard}

        \paragraph{Orthodox motivation.}
        \textbf{[ORTHODOX INPUT].}
        Maxwell's kinetic-theory analysis separates collective translation from
        internal rotational motion and shows that, when additional internal
        degrees of freedom are dynamically coupled and equilibrate, they alter
        the thermodynamic energy bookkeeping and therefore the predicted heat
        capacity \cite{Maxwell1860GasesI,Maxwell1860GasesIIIII}.

        \paragraph{SST consistency consequence.}
        \textbf{[DERIVED / CONSISTENCY GATE].}
        For a candidate SST knot class $K$, an internal mode $a$ is
        phenomenologically active only if all three conditions are satisfied:
        \begin{align}
            \mathcal G_{K,a} &\neq 0,
            \label{eq:canon_mode_coupling_nonzero}\\
            E_{\rm drive} &\gtrsim \Delta_{K,a},
            \label{eq:canon_mode_kinematic_access}\\
            \tau_{K,a} &\lesssim t_{\rm obs},
            \label{eq:canon_mode_equilibration_time}
        \end{align}
        where $\mathcal G_{K,a}$ denotes the interaction-to-mode coupling
        amplitude, $\Delta_{K,a}$ the true excitation-energy threshold, and
        $\tau_{K,a}$ the excitation/equilibration timescale.  Failure of any one
        condition can freeze the mode on the relevant experimental timescale.

        \paragraph{Frequency--gap distinction.}
        \textbf{[ORTHODOX INPUT / CANON GUARD].}
        A nonzero classical normal-mode frequency does not by itself imply a
        nonzero excitation-energy gap:
        \begin{equation}
            \omega_{K,a}>0
            \not\Rightarrow
            \Delta_{K,a}>0 .
            \label{eq:canon_frequency_not_gap}
        \end{equation}
        A classical harmonic mode with energy quadratic in its amplitude admits
        arbitrarily small excitation energy.  Therefore a positive SST gap must
        arise from an independently demonstrated discrete occupancy rule,
        topological barrier, nonlinear threshold, or other mechanism that
        forbids arbitrarily small excitation energies.  Stiffness alone is not
        sufficient.

        \paragraph{Mode-count gate.}
        \textbf{[DERIVED / CONDITIONAL].}
        If several ungapped internal modes are appreciably coupled and
        equilibrate on the measurement timescale, orthodox statistical mechanics
        assigns them an order-$k_B$ thermodynamic contribution per active
        quadratic degree of freedom.  Consequently, a particle assignment that
        predicts many such modes but no corresponding heat-capacity, inelastic,
        broadening, or sideband signatures fails this gate unless the coupling or
        equilibration assumptions are independently shown to be invalid
        \cite{PathriaBeale2011}.

        \paragraph{Research-track pointer.}
        \textbf{[SPECULATIVE / RESEARCH PROGRAMME].}
        The companion research track develops the full
        Maxwell--SST kinetic closure: phase-space distributions for knot
        ensembles, interaction-induced mode-coupling matrices, a
        translation/orientation/Kelvin/twist/writhe/core classification,
        a separate knot-ensemble stress tensor, and an isotropization test on
        $SO(3)$.

% ============================================================
% Placement B:
% Replace the existing \subsection{Canonical Statement}
% with the block below.
% ============================================================
    \subsection{Canonical Statement}
        \label{subsec:canon_internal_mode_statement}

        \textbf{[DERIVED / CONSISTENCY PRINCIPLE].}
        \begin{center}
            \begin{minipage}{0.92\linewidth}
                \itshape
                Any viable SST knot assignment must account for every physically
                coupled internal mode.  Each such mode must be shown either to
                possess a true excitation-energy gap sufficient to suppress its
                occupation, to couple weakly enough to satisfy empirical bounds, or
                to remain dynamically unequilibrated on the relevant observation
                timescale.  The same preregistered mode ledger must satisfy
                spectroscopy and thermodynamics without post-hoc deletion of
                coupled degrees of freedom.
            \end{minipage}
        \end{center}
